
\documentclass[11pt, a4paper]{article}

% ===== ESSENTIAL PACKAGES =====
\usepackage{geometry}
\geometry{a4paper, margin=1.2in}
\usepackage{setspace}
\usepackage{array}
\usepackage{xcolor}
\onehalfspacing

\usepackage{catchfile}

% ===== WORD COUNT COMMAND =====
\newcommand{\wordcount}{%
    \ifnum\pdf@shellescape=1%
        \immediate\write18{texcount -1 -sum Oblig1.tex > wordcount.txt}%
        \CatchFileDef{\words}{wordcount.txt}{}%
        \words%
    \else%
        [Enable -shell-escape]%
    \fi%
}
\usepackage{mdframed}

\newmdenv[
    leftmargin=-40pt,
    rightmargin=0pt,
    innerleftmargin=3pt,
    innerrightmargin=0pt,
    innertopmargin=5pt,
    innerbottommargin=5pt,
    skipabove=5pt,
    skipbelow=5pt,
    linewidth=1pt,
    linecolor=black,
    topline=true,
    rightline=false,
    bottomline=false
]{sectionboxinner}

\newenvironment{sectionbox}[1][]{%
    \begin{sectionboxinner}
    \textbf{\large #1}
    \vspace{1pt}
    \newline
}{%
    \end{sectionboxinner}
}

\newmdenv[
    leftmargin=0pt,
    rightmargin=0pt,
    innerleftmargin=0pt,
    innerrightmargin=0pt,
    innertopmargin=3pt,
    innerbottommargin=3pt,
    skipabove=8pt,
    skipbelow=8pt,
    linewidth=1.2pt,
    linecolor=gray!60,  % Lighter gray color
    topline=false,
    rightline=false,
    bottomline=false
]{subsectionboxinner}

\newenvironment{subsectionbox}[1][]{%
    \begin{subsectionboxinner}
    \textbf{#1}
    \vspace{3pt}
    \newline
}{%
    \end{subsectionboxinner}
}

% ===== CUSTOM ENVIRONMENTS FOR PHILOSOPHICAL ARGUMENTS =====
\newenvironment{argument}{\begin{quote}\itshape\textbf{Argument: }}{\end{quote}}
\newenvironment{objection}{\begin{quote}\itshape\textbf{Innvending: }}{\end{quote}}
\newenvironment{reply}{\begin{quote}\itshape\textbf{Reply: }}{\end{quote}}

% ===== DOCUMENT BEGIN =====
\title{Kap. 11 - Evolusjon. Er vi overlevelsesmaskiner skapt av genene våre?}
\author{Oliver Ekeberg}
\date{\today}



\begin{document}
\maketitle

\tableofcontents

\begin{sectionbox}[Den beste ideen noensinne?]

\begin{flushleft}
    Filosofen Daniel Bennet mener at den beste ideen noensinne går til Charles Darwins teori om evolisjon. For å være presis, skiller teorien om evolusjon gjennom naturlig utvalg, fordi den forener to områder; liv, mening og formål med rom og tid, årsak og virkning, mekanismer og fysiske lover.
\end{flushleft}

\begin{flushleft}
    Biologen \textit{Richard Dawkins} sa han er enig med Bennet. Han sa at hvis utenomjordisk liv besøker oss, vil det første de spør være, om menneskene har funnet en form for evolusjonsteori. Vi skal her utforske den grunnleggende strukturelle teorien om evolusjon gjennom naturlig utvalg. og Dawkins ideer om hva evolusjonsteorien innebærer for menneskets plass i naturen
\end{flushleft}
\end{sectionbox}

\begin{sectionbox}[Strukturen til evolusjon gjennom naturlig utvalg]
\begin{flushleft}
    Vi skal nå se evolusjon gjennom naturlig utvalg. 
    \begin{quote}
        \textbf{\textit{Evolusjon gjennom naturlig utvag}}: En prosess der organisme populasjoner endres over tid. En populasjon er en gruppe, vanligvis av samme art. For eksempel svarttroster i Norge\\

        Denne "prosessen" har tre forutsetninger
        \begin{enumerate}
            \item Variasjon
            \item Arv
            \item Fitness forskjeller
        \end{enumerate}
    \end{quote}
\end{flushleft}

\begin{flushleft}
    \begin{enumerate}
        \item \textbf{\textit{Variasjon:}} Dette vil si at en gruppe organisme har variasjoner i karakteristikker, de skiller seg fra hverandre i morfologi, fysiologi og oppførsel. En svarttrost populasjon kan ha forskjellig lengde, annerledes nebb osv.
        \item \textbf{\textit{Arv:}} Variasjon er delvis arvelig. dvs. avkom ligner mer på sine foreldre enn ubeslektede individer. 
        \item \textbf{\textit{Fitness-forskjeller:}} De forskjellige variantene i en populasjon, etterlater seg forskjellige antall avkom. Si at svarttroster med lengre vinger får flere unger. Da vil også ungene deres ha lengre vinger. Over tid, vil gjennomsnittet konvergere mot et nytt likevektspunkt. Denne konvergensen kan beskrives som at populasjonen utvikler seg\\
        

        Når du hører fitness-forskjeller, tenk "gjennomsnittlig antall avkom over tid"
    \end{enumerate}
\end{flushleft}

\begin{subsectionbox}[Evolusjon og befolkningsvekst]
    Det er også viktig å regne med at mange svarttroster dør ut før de kan få barn. Dette er inspirert av økonomen Thomas Malthus
\end{subsectionbox}

De tre forutsetngene (variasjon, arv og fitness) er kjernen i teorien om evolusjon gjennom naturlig utvalg. Så lenge en populasjon har de tre forutsetningene, så vil den populasjonen utvikle seg. Modellen er en svært enkel og systematisk prosess. Dennett og Dawkins er blitt imponert over hvor mye som kan forklares av den, straks den anvendes på konkrete tilfeller.
\end{sectionbox}

\begin{sectionbox}[Gener]
    \begin{flushleft}
        Da Darwin kom med evolusjonsteorien, var det et stort hull mellom korrelasjon og kausalitet. Hva var det som fikk avkom av foreldre til å likne mer på de, enn på ubeslektede? Det kom ikke en pekepinn på en forklaring før introduksjonen av DNA kom inn på 1950-tallet. Men tidligere, var det Gregor Mendel på 1900-tallet som mente at biologisk arv hadde å gjøre med \textbf{\textit{separate enheter}}, disse enhetene ble kalt for \textbf{\textit{gener}}. Et DNA er en streng av gener, som blir kopiert fra foreldrene til avkommet. Denne kopierende strukturen til DNA hjelper til med å tette gapet mellom korrelasjon og kausalitet i evolusjonsteori.\\
    \end{flushleft}

    \begin{flushleft}
        Ifølge Dawkins begynner evolusjonen med "replikatorer". Gener er de kraftigste replikatorer
    \end{flushleft}
\end{sectionbox}

\begin{sectionbox}[Tilfeldighet]
    \begin{flushleft}
        Vi har forklart at evolusjon ved naturlig seleksjon har tre forutsetninger; variasjon, arv og fitness-forskjeller. Teorien sier at \textbf{\textit{variasjonene er arvelige og leder til at fitness-forskjeller akkumulerer over tid}}. \\

        Hva er da oppgavet til variasjonene? Noen variasjon kommer fra trening og læring. For eksempel at en svarttrost lærer å fly en manøvrer av foreldrene sine. Slik tilegnet variasjon er stort sett ikke arvelig. Poenget er at det er ingen pålitelig mekanisme som vil "overføre" det tilegnede trekket til neste generasjon. Siden tilegnet variasjon \textit{ikke} er arvelig, vil den ikke lede til evolusjonær endring. \\

        Hvor kommer da arvelig variasjon fra? For Dawkins gensentrerte syn er tilfeldighet svaret. Du kan si at en følge er konvergent, selvom den har tilfeldig støy i seg. Hver gang DNA fra to individer blander seg, vil det lage en ny kombiasjon med gener, og gjøre det nye individet unikt.
    \end{flushleft}
\end{sectionbox}

\begin{sectionbox}[Egoisme]
    \begin{flushleft}
        På hvilken måte er gener egoistiske? Et gen er et molekyl, og er ikke noe som vil ha noe. Fremdeles argumenterer Dawkins at det er som om gener bare bryr om seg selv. \\
        
        \begin{quote}
            Dawkin: Han argumenterer for at mennesker og andre organisme ikke er mer enn maskiner skapt av gener for deres egen evolusjonære suksess
        \end{quote}
        Genene som lager flere kopier av seg selv, enn andre vil vinne evolusjonen. Dette sikres av prosessen for evolusjon gjennom naturlig utvalg. Hvis mennekser ble til gjennom denne prosesse, må det ha vært fordi vi har hatt gener som er flinkere til å replikere oss selv enn våre konkurrenter. Det er derfor Neanderthalene for eksempel er en utdødd menneskeart. \\

        Dawkins gender at gener er egoistiske på den måten at den grunnleggende forklaringen på alle trekkene våre er at eksistensen til disse hjalp noen gener til å bli bedre til å replikere seg selv enn konkurrentene deres. Dette er det eneste som har noe å si for forklaringen av alle trekkene våre.
    \end{flushleft}
\end{sectionbox}

\begin{sectionbox}[Studiespørsmål]

    \begin{enumerate}
        \item Hva er replikatorer, ifølge Dawkins? Hva er Dawkins grunn til å tro at de er så viktige?
        \item Mange mennesker på står at teorien om evolusjon gjennom naturlig utvalg er det endelige beviset mot det teologiske verdensbildet vi støtter på hos Aristoteles. Evolusjon forklarer hvorfor ting i naturen ser ut til å ha en hensikt, selv om de faktisk er et resultat av tilfeldighet. Hva er en god grunn for å to dette om evolusjon, er det riktig?
        \item Lykkes Dawkins i å vise at den grunnleggende forklaringen på alle trekkene våre er at de hjalp noen gener med å replikere seg selv?
        \item Hva, om noe, sier teorien om evolusjon gjennom naturlig utvalg om hvorvidt vi mennesker er naturlig egoistiske eller ikke?
    \end{enumerate}

    \textbf{\textit{Dette bør du kunne}}
    \begin{itemize}
        \item Forstå den grunnleggende strukturen til teorien om evolusjon gjennom naturlig uvalg og være i stand til å illustrere den med et eksempel
        \item Forstå genenes rolle i Dawkins syn på evolusjon gjennom naturlig utvalg
        \item Forstå hvilken måte Dawkins mener at gener er egoistitske, og være i stand til å reflektere kritisk over hva dette innebærer for menneskers natur
    \end{itemize}
\end{sectionbox}

\end{document}